\begin{afterword}
\summaryhead

The post opens with a choice: save 400 lives for certain, or save 500 with probability 90\%. Most people, it says, take the sure option, which is foolish, because the gamble saves 450 lives on average. Described as deaths (100 die for certain, or a 10\% chance that all 500 die), the same options lead a majority to take the gamble. The preference follows the feeling that the wording produces, and a life is worth more than that feeling. Hence the rule: ``shut up and multiply.''

The same rule means that if a hiccup has any disvalue, some number of hiccups, a googolplex for instance, is worse than one person's death by shark attack. Such trade-offs are everyday choices, such as a console game against charity, or one bone marrow transplant against 200 children saved from malaria. People avoid them by calling lives ``sacred'' and taking offence at any trade with money. The post ends by saying that altruism is judged by the help it gives, not by how it feels, and that rationality rewards those who give themselves to it.

\responsehead

The centre of the post is sound, and it is Tversky and Kahneman's. In their 1981 study the same pair of options, described as lives saved and then as lives lost, drew opposite majorities. Whatever the right choice is, a preference that reverses with the wording is following the words, not the lives.

Its own numbers carry a claim it does not source. In the original problem the two options had the same expected value, so the majority who took the sure option were not choosing the worse option on average. The post raises the gamble's expected value to 450 against 400, says that ``Most people'' still choose the sure option, and calls them foolish. No study is cited for those numbers.

The second case asks more than the post supplies. It says that if hiccups have ``any'' disvalue, then ``on pain of decision-theoretic incoherence'' enough of them must outweigh a shark attack. Coherence alone does not require this: a ranking in which each extra hiccup is worse but the total never reaches the shark attack is consistent. What the conclusion needs is that each hiccup adds the same amount however many there are. Harsanyi's theorem gives the best-known argument for that premise, and the next post in the book gives the author's. The first version of this post, ``Circular Altruism'', also argued for the conclusion, with a chain of small steps from torture to dust specks. The book's version keeps the conclusion and drops the chain, so here the claim stands on assertion.

The evidence on sacred values is Philip Tetlock's and is fairly reported. The anecdote about the government agency has no source, and the author's first version said it had none to hand. The closing promise that rationality ``gives to you in return'' is exhortation, with no evidence offered.

In short: a correct and well-chosen demonstration of the framing effect, taken from Tversky and Kahneman, followed by a claim about hiccups and shark attacks that coherence alone does not establish and that this version asserts without the argument its first version gave.
\end{afterword}
